% =====================================================================
%  KONCEPT KAPITOLY: METODIKA VLASTNÉHO VÝSKUMU  (draft pre DP)
%  ---------------------------------------------------------------------
%  ⚠️  POZOR — PRAVIDLÁ PRÁCE:
%  Toto je AI-ASISTOVANÝ DRAFT. Podľa zadania DP musí byť AI-generovaný
%  obsah (a) označený vertikálnou čiarou a (b) priznaný v doplnkovom
%  materiáli. Tento súbor preto NEVKLADAJ priamo bez revízie. Postup:
%    1) Prečítaj, preformuluj vlastnými slovami a over každé tvrdenie.
%    2) Časti, ktoré ostanú AI-asistované, nechaj v prostredí {aibar}.
%    3) Pridaj zdroj do doplnkového materiálu (priznanie AI).
%  Štruktúra a tvrdenia vychádzajú z podkladov 01–05 v tomto priečinku
%  a z tvojej existujúcej SOTA kapitoly.
%
%  Vyžaduje v preambule:  \usepackage{framed}  (pre vertikálnu čiaru).
%  Ak framed nemáš, definovaný je fallback (obyčajný odsek).
% =====================================================================

% --- značenie AI-asistovaného obsahu vertikálnou čiarou ---
\ifcsname leftbar\endcsname
  \newenvironment{aibar}{\begin{leftbar}}{\end{leftbar}}
\else
  \newenvironment{aibar}{\par\noindent\ignorespaces}{\par}
\fi

\chapter{Metodika vlastného výskumu}
\label{chapter:methodology}

\begin{aibar}

Táto kapitola transformuje východiská zo súčasného stavu poznania
(kapitola~\ref{chapter:sota}) do konkrétneho, realizovateľného výskumného
dizajnu. Nadväzuje na záver rešerše, podľa ktorého je pre slovenský kontext
najvhodnejšie kombinovať tri prístupy: trojskupinový dizajn (Slijepcevic a
Yaylali), systematický zber interakčných dát (Lee a kolektív) a doplnenie
self-report meraní o behaviorálne ukazovatele (Kosmyna a kolektív).


\section{Výskumné otázky a hypotézy}
\label{sec:rq}

Výskum sleduje tri vzájomne prepojené výskumné otázky, ktoré zodpovedajú
trom medzeram identifikovaným v rešerši (dominancia self-reportu, hrubá
typológia používateľov, podreprezentovanosť slovenského kontextu).

\paragraph{RQ1 (deskriptívna).}
Ako slovenskí stredoškoláci reálne využívajú generické a adaptívne nástroje
umelej inteligencie? Predpokladáme existenciu aspoň troch odlíšiteľných
profilov používateľov (H1a), pričom profil sa líši podľa typu školy (H1b).

\paragraph{RQ2 (experimentálna).}
Aký vplyv má typ použitého nástroja na vedomostný zisk a na kritické zapojenie
žiaka? Predpokladáme vyšší vedomostný zisk pri adaptívnom nástroji (H2a) a
nižšie kritické zapojenie pri generickom nástroji, merané behaviorálnymi
ukazovateľmi (H2b).

\paragraph{RQ3 (vysvetľujúca).}
Akým mechanizmom ovplyvňuje používanie umelej inteligencie kritické myslenie?
Predpokladáme, že vzťah je mediovaný kognitívnym offloadingom (H3a) a že
vyššia dôvera v nástroj súvisí s nižším kritickým myslením, zatiaľ čo vyššia
sebadôvera s vyšším (H3b).


\section{Výskumný dizajn}
\label{sec:design}

Zvolený je trojskupinový kvázi-experiment s meraním pred intervenciou a po
nej. Účastníci sú priradení do troch podmienok: (A)~adaptívny nástroj,
(B)~generický konverzačný nástroj a (C)~kontrolná skupina bez asistencie
umelej inteligencie. Meracie nástroje pre-testu a post-testu sú ekvivalentné,
no nie identické, čím sa eliminuje efekt opakovaného testovania.

Vzhľadom na to, že historická vzorka žiakov z obdobia pred rozšírením
generatívnej umelej inteligencie na Slovensku neexistuje a porovnanie
nesúrodých kohort by bolo metodologicky neinterpretovateľné, je úloha
``pred-AI'' referencie nahradená synchrónnou kontrolnou skupinou~(C).
Doplnkovo je dizajn rozšírený o longitudinálne jadro s odloženým re-testom,
ktoré umožňuje sledovať retenciu a aproximovať fenomén kognitívneho dlhu
dokumentovaný neurofyziologicky v štúdii Kosmyny a kolektívu.


\section{Vzorka a výber škôl}
\label{sec:sample}

Jadrom vzorky je gymnázium (2.--3. ročník), ktoré predstavuje najhomogénnejšiu
a najľahšie interpretovateľnú populáciu s dostupnými štandardizovanými testami.
Vzorka je rozšírená o jednu technickú strednú odbornú školu ako kontrastnú
podmienku, čo umožňuje testovať hypotézu H1b o rozdieloch medzi typmi škôl.
Cieľová veľkosť je orientačne 30--60 účastníkov na typ školy; analýza sily pri
troch skupinách, strednom efekte, hladine $\alpha = 0{,}05$ a sile $0{,}8$
stanovuje minimum na 18~účastníkov na skupinu, k čomu je pripočítaná rezerva na
výpadky.


\section{Meracie nástroje}
\label{sec:instruments}

Meranie je organizované do štyroch vrstiev (tabuľka~\ref{tab:vrstvy}).
Doménový vedomostný zisk je meraný adaptovaným konceptuálnym testom
(\emph{concept inventory}). Kritické myslenie je merané vlastnou batériou úloh
mapovanou na päť dimenzií Halpernovho modelu, doplnenou o Test kognitívnej
reflexie; Halpernov komerčný test (HCTA) je pritom použitý len ako teoretický
rámec, nie ako licencovaný nástroj. Vzťah k umelej inteligencii je meraný
voľne dostupnými škálami (kognitívne offloading, dôvera v nástroj, gramotnosť v
oblasti umelej inteligencie). Štvrtou, originálnou vrstvou sú behaviorálne
ukazovatele extrahované z interakčných logov, ktoré slúžia ako dostupná
náhrada za neurofyziologické meranie. Všeobecná rozumová schopnosť
(Ravenove matice, prípadne voľne dostupný itemset ICAR) je meraná len ako
kontrolná premenná.

\begin{table}[H]
\centering
\caption{Štyri vrstvy merania}
\label{tab:vrstvy}
\small
\begin{tabularx}{\textwidth}{|l|X|X|}
\hline
\textbf{Vrstva} & \textbf{Konštrukt} & \textbf{Nástroj} \\
\hline
A & Vedomostný zisk & Adaptovaný concept inventory \\
\hline
B & Kritické myslenie, metakognícia & Vlastná batéria (5 dimenzií) + CRT + MAI \\
\hline
C & Vzťah k umelej inteligencii & Offloading, dôvera (SAID), AI gramotnosť \\
\hline
D & Behaviorálne zapojenie & Logy: čas čítania, miera úprav, doplňujúce dopyty \\
\hline
\end{tabularx}
\end{table}

Pred ostrým zberom je naplánovaná pilotná validácia vlastnej batérie a hernej
aplikácie na vzorke približne 30~žiakov (položková analýza, vnútorná
konzistencia).


\section{Nástroj zberu dát a herná aplikácia}
\label{sec:tool}

Zber prebieha dvoma komplementárnymi kanálmi. Prvým je sieťové zachytávanie
reálnych dopytov prostredníctvom lokálneho proxy (popísané v
kapitole~\ref{chapter:sota} a v praktickej časti), ktoré poskytuje vysokú
ekologickú validitu. Druhým je herná aplikácia ``AI~Detektív'', ktorá
prevádza schopnosti žiaka na uzavreté voľby a tým rieši problém obtiažnej
kvantifikácie otvorených odpovedí. V hre žiak rieši úlohy s asistentom, ktorý
zámerne v časti prípadov poskytuje nesprávne informácie; pri každom
rozhodovacom uzle žiak volí z uzavretých možností, ktoré sú napevno priradené
k meranej dimenzii a skórované automaticky. Tri režimy hry (adaptívny,
generický, bez asistencie) zodpovedajú trom experimentálnym skupinám, pričom
meracie uzly sú vo všetkých režimoch identické, čím je zachovaná
porovnateľnosť.


\section{Analýza dát}
\label{sec:analysis}

Vplyv intervencie na vedomostný zisk je analyzovaný kovariančnou analýzou
(ANCOVA) s post-testom ako závislou premennou, skupinou ako faktorom a
pre-testom ako kovariátom; doplnkovo sú uvádzané normalizovaný zisk a
Cohenovo~$d$ kvôli porovnateľnosti so zahraničnými štúdiami. Zmena v čase je
modelovaná lineárnym zmiešaným modelom s náhodným interceptom na účastníka.
Mechanizmus vplyvu na kritické myslenie je testovaný mediačnou analýzou s
bootstrapovaným intervalom spoľahlivosti nepriameho efektu. Profily
používateľov sú identifikované zhlukovou analýzou nad behaviorálnymi a
self-report premennými. Kvalitatívne dáta z pološtruktúrovaných rozhovorov sú
spracované tematickou analýzou a slúžia na interpretáciu kvantitatívnych
nálezov.


\section{Etické aspekty a ochrana údajov}
\label{sec:ethics}

Keďže účastníkmi sú prevažne maloletí žiaci, výskum si vyžaduje informovaný
súhlas zákonných zástupcov aj samotných žiakov. Pred zberom je vykonaná analýza
vplyvu na ochranu údajov (DPIA) podľa vzoru fínskeho modelu a získané je
schválenie etickej komisie fakulty. Dáta sú pseudonymizované už pri zázname;
surové logy sú uchovávané len po dobu nevyhnutnú na analýzu a následne
nahradené agregátmi.

\end{aibar}
